\begingroup
\makeatletter\tagsleft@true\let\veqno\@@leqno\makeatother
\newcommand{\TransBracket}[1]{\boldsymbol{[}#1\boldsymbol{]}}
\ArithmeticHead{296}{L}{ANALYSIS. NACHLASS.}
\NumberedItem{9}
\par\indent Durch Induction gefunden:
\[\varepsilon-\varepsilon\varepsilon=\surd3\]
\[\begin{aligned}
&\{1+2x+2x^4+2x^9+{.}{.}{.}\}+(\varepsilon-\varepsilon\varepsilon)\{1+2x^3+2x^{12}+2x^{27}+{.}{.}{.}\}\\
&=-2\varepsilon\varepsilon(1-\varepsilon x)(1-\varepsilon\varepsilon x^3)(1-\varepsilon x^5)(1-\varepsilon\varepsilon x^7){.}{.}{.}\\
&\qquad(1+x^3)(1+x^9)(1+x^{15})(1+x^{21}){.}{.}{.}\\
&\qquad\times(1-x^4)(1-x^8)(1-x^{12}){.}{.}{.}
\end{aligned}\]
\noindent Dieselbe Summe wäre demnach
\[\begin{aligned}
&=-2\varepsilon\varepsilon\{1-\varepsilon x-\varepsilon\varepsilon x^3+\varepsilon\varepsilon x^6+\varepsilon x^{10}-x^{15}-x^{21}+\varepsilon x^{28}+\varepsilon\varepsilon x^{36}-\text{etc.}\}\\
&\qquad.(1+x^3)(1+x^9)(1+x^{15}){.}{.}{.}\\
&=\frac2{x^{\frac18}}\{-\varepsilon\varepsilon x^{\frac18}+x^{\frac98}+\varepsilon x^{\frac{25}8}-\varepsilon x^{\frac{49}8}-x^{\frac{81}8}+\varepsilon\varepsilon x^{\frac{121}8}-\text{etc.}\}\\
&\qquad(1+x^3)(1+x^9)(1+x^{15})\ \text{etc.}
\end{aligned}\]
\clearpage
\ArithmeticHead{297}{R}{{\scriptsize ZUR THEORIE DER TRANSSCENDENTEN FUNCTIONEN. I}}
\SourceLeaf{304}
\noindent Woraus folgt
\begin{equation*}\tag*{1.}
\begin{aligned}
&1+2x+2x^4+2x^9+\text{etc.}\\
&=\frac1{x^{\frac18}}\{x^{\frac18}+2x^{\frac98}-x^{\frac{25}8}+x^{\frac{49}8}-2x^{\frac{81}8}-{.}{.}{.}\}\\
&\qquad.(1+x^3)(1+x^9)(1+x^{15}){.}{.}{.}
\end{aligned}
\end{equation*}
\begin{equation*}\tag*{2.}
\begin{aligned}
&1+2x^3+2x^{12}+2x^{27}+\text{etc.}\\
&=\frac1{x^{\frac18}}\{x^{\frac18}+x^{\frac{25}8}-x^{\frac{49}8}-x^{\frac{121}8}+x^{\frac{169}8}+{.}{.}{.}\}\\
&\qquad.(1+x^3)(1+x^9)(1+x^{15}){.}{.}{.}[\text{*)}]
\end{aligned}
\end{equation*}
\[\text{Factor communis}=\frac{\TransBracket{x^6}^2}{\TransBracket{x^3}\TransBracket{x^{12}}}\mathpunct{\raisebox{.5ex}{.}}\]
\par\indent Auch scheint zu seyn
\begin{equation*}\tag*{[9]}
\left\{\begin{aligned}
&1+2x^9+2x^{36}+\text{etc.}\\
&=\frac1{x^{\frac18}}\{x^{\frac18}-x^{\frac{25}8}+x^{\frac{49}8}-x^{\frac{121}8}-x^{\frac{169}8}{.}{.}{.}\}\\
&\qquad(1+x^3)(1+x^9)(1+x^{15}){.}{.}{.}
\end{aligned}\right.
\end{equation*}
\begin{equation*}\tag*{[10]}
\left\{\begin{aligned}
&x+x^4+x^{16}+x^{25}+\text{etc.}\\
&=\frac1{x^{\frac18}}\{x^{\frac98}-x^{\frac{81}8}-x^{\frac{225}8}+x^{\frac{441}8}+x^{\frac{729}8}-{.}{.}{.}\}\\
&\qquad(1+x^3)(1+x^9)(1+x^{15}){.}{.}{.}\\
&=x\frac{\TransBracket{x^6}^2\TransBracket{x^9}\TransBracket{x^{36}}}{\TransBracket{x^3}\TransBracket{x^{12}}\TransBracket{x^{18}}},
\end{aligned}\right.
\end{equation*}
\noindent folglich
\[1-x-x^3+x^6+\text{etc.}=\frac{\TransBracket{x}\TransBracket{x^4}}{\TransBracket{xx}},\]
\noindent wie gehörig.
\begin{center}\textit{Ist alles bewiesen.}\end{center}
\Needspace{3\baselineskip}
\par\smallskip\noindent\rule{.25\linewidth}{.4pt}\par
{\small\par\indent [*) In der Handschrift steht $-x^{\frac{169}8}-{.}{.}{.}$]\par}
\clearpage
\endgroup
